% =============================================================================
\section{Tarefa 3 --- Execução no robô}
\label{sec:exec}
% =============================================================================

\subsection{Script}

O script \texttt{src/run\_robot.py} usa só funções do SDK~\cite{xarm_sdk}:
gera a trajetória e corre a verificação da Secção~\ref{sec:viab} (se alguma
condição falhar, nada é enviado); liga ao robô, lê o desvio de ferramenta e
converte as poses; aproxima-se da pose inicial junta a junta, a
\SI{20}{\degree\per\second}; e executa o swing por um de dois métodos. No
método \texttt{servo} (modo~1) envia uma pose a cada \SI{10}{\milli\second}
com \texttt{set\_servo\_cartesian} (Listagem~\ref{lst:servo}): o controlador
resolve a inversa em cada ponto e o perfil temporal é o nosso, incluindo a
aceleração constante do downswing. No método \texttt{pontos} (modo~0) envia um
ponto em cada dez com \texttt{set\_position} e o controlador interpola entre
eles, impondo o seu próprio perfil de velocidade. O parâmetro
\texttt{-{}-lento N} estica o tempo $N$ vezes, mantendo o caminho.

\begin{lstlisting}[language=Python, caption={Ciclo do método \texttt{servo}
(simplificado de \texttt{src/run\_robot.py}).},
label={lst:servo}, captionpos=b]
arm.set_mode(1); arm.set_state(0)            # modo servo
t0 = time.perf_counter()
for k, pose in enumerate(poses):             # [x, y, z, roll, pitch, yaw]
    if arm.set_servo_cartesian(pose) != 0:   # erro do controlador: parar
        break
    _, medida = arm.get_position()
    registo.append([time.perf_counter() - t0, *pose, *medida])
    while time.perf_counter() - t0 < (k + 1) * dt:
        pass                                 # espera até ao instante seguinte
arm.set_mode(0); arm.set_state(0)
\end{lstlisting}

\subsection{Resultados}

Os ensaios da versão final fizeram-se a 6/11/2026 (\emph{firmware} 2.7.1). O
erro é a distância entre o flange medido e o pedido no mesmo instante; o
atraso é o deslocamento temporal que minimiza esse erro
(Tabela~\ref{tab:exec}, Figura~\ref{fig:exec}). O método \texttt{pontos}
também executou o swing completo, com \texttt{-{}-lento 3}, mas não regista a
pose medida.

\begin{table}[H]
\centering
\caption{Seguimento em modo \texttt{servo} (flange); cada ensaio foi feito
duas vezes e os dois dão o mesmo erro.}
\label{tab:exec}
\small
\begin{tabular}{lcccccc}
\toprule
\rowcolor{tabcab}
Ensaio & Duração & Pontos & Erro médio & Erro máx. & Atraso & Erro sem o atraso \\
\midrule
\texttt{-{}-lento 2}, 1.º & \SI{10.8}{\second} & 1085/1085 & \SI{0.54}{\milli\meter} & \SI{2.5}{\milli\meter} & \SI{10}{\milli\second} & \SI{0.17}{\milli\meter} \\
\texttt{-{}-lento 2}, 2.º & \SI{10.9}{\second} & 1085/1085 & \SI{0.57}{\milli\meter} & \SI{3.2}{\milli\meter} & \SI{10}{\milli\second} & \SI{0.27}{\milli\meter} \\
\texttt{-{}-lento 1}, 1.º & \SI{5.4}{\second}  & 543/543   & \SI{8.4}{\milli\meter}  & \SI{28.0}{\milli\meter} & \SI{58}{\milli\second} & \SI{7.2}{\milli\meter} \\
\texttt{-{}-lento 1}, 2.º & \SI{5.4}{\second}  & 543/543   & \SI{8.4}{\milli\meter}  & \SI{28.2}{\milli\meter} & \SI{55}{\milli\second} & \SI{7.2}{\milli\meter} \\
\bottomrule
\end{tabular}
\end{table}

\begin{figure}[H]
\centering
\fig{Photos/run_servo.png}{0.92\textwidth}
\caption{Posição do flange pedida (tracejado) e medida, e erro de posição, a
metade da velocidade (esquerda) e à velocidade planeada (direita). As linhas
verticais separam as três fases.}
\label{fig:exec}
\end{figure}

\begin{itemize}[itemsep=3pt]
  \item \textbf{A metade da velocidade o robô segue a trajetória.} O erro
        médio é de \SI{0.5}{\milli\meter}, da ordem da repetibilidade do robô
        ($\pm\SI{0.5}{\milli\meter}$~\cite{ufactory_lite6}), e quase todo é
        um atraso de um ciclo (\SI{10}{\milli\second}).
  \item \textbf{À velocidade planeada o robô atrasa-se onde a velocidade
        varia depressa}: no backswing (média \SI{11}{\milli\meter}) e no
        follow-through (\SI{16}{\milli\meter}); no downswing, de aceleração
        constante, o erro é de \SI{1.4}{\milli\meter}. Não é um atraso puro
        (descontar \SI{58}{\milli\second} só baixa o erro para
        \SI{7.2}{\milli\meter}), nem saturação de velocidade (a máxima
        medida, \SI{173}{\milli\meter\per\second}, iguala a pedida), nem o
        limite de aceleração das juntas (a trajetória pede no máximo
        \SI{227}{\degree\per\second\squared}, um quinto do do fabricante).
        É compatível com um limite de aceleração cartesiana ou com
        uma filtragem do controlador em modo servo, que sem acesso à sua
        configuração não foi possível distinguir.
        
\end{itemize}

\subsection{Ocorrências e correções}

Na primeira sessão (5/11) o follow-through ia até $\varphi_{end} = \ang{70}$
e a pinça ainda girava sobre si própria. Duas ocorrências levaram à versão
final:

\begin{itemize}[itemsep=3pt]
  \item \textbf{Erro C31 (deteção de colisão).} Com $\varphi_{end}=\ang{70}$
        o robô parou sempre no mesmo ponto, quase parado e com o braço
        esticado para o lado (junta~5 a \ang{110}). Não sendo um efeito
        dinâmico, as hipóteses são o cabo da pinça a esticar ou a massa da
        pinça não estar configurada no controlador. Encurtar o follow-through
        para \ang{50} eliminou o erro.
        
  \item \textbf{Erro C24 no método \texttt{pontos}.} À velocidade
        planeada, a interpolação que o próprio controlador faz entre os pontos
        enviados pedia às juntas mais velocidade do que a trajetória, e o
        controlador interrompeu o movimento, como lhe compete. A trajetória
        planeada respeita os limites (Tabela~\ref{tab:viab}); por isso o
        método \texttt{pontos} só se usa a velocidade reduzida ($N \geq 2$).
\end{itemize}

A orientação final (Secção~\ref{sec:traj}) reduziu o erro de seguimento de
2,0 para \SI{0.5}{\milli\meter} a metade da velocidade e de 18,9 para
\SI{8.4}{\milli\meter} à velocidade planeada. A verificação de viabilidade
passou também a incluir a velocidade do flange, que é o ponto comandado.
